\section{Exact inverse structure}
\label{sec:inverse}

\subsection{Triangular inverse}

\begin{theorem}[Conditional exact inverse]
\label{thm:inverse}
Fix a wall-free one-step advance \cref{eq:advance}. If \(\gamma\neq0\) and the phase map \(G_{\bm r^+}\) is bijective at the advanced position, then the complete pre-state is uniquely recovered by
\begin{subequations}
\label{eq:inverse}
\begin{align}
\bm\phi
  &=G_{\bm r^+}^{-1}(\bm\phi^+),
  \label{eq:inverse-phase}\\
\bm v
  &=\gamma^{-1}
  \left[
  \bm v^+
  -\dt M^{-1}
  \bm F(\bm r^+,\bm\theta,\bm\phi)
  \right],
  \label{eq:inverse-velocity}\\
\bm r
  &=\bm r^+-\bm v\dt,
  \label{eq:inverse-position}\\
\bm\theta&=\bm\theta^+.
\end{align}
\end{subequations}
\end{theorem}

\begin{proof}
The final position is known. Bijectivity of \(G_{\bm r^+}\) determines the old phase. The force is then known because it depends on \((\bm r^+,\bm\theta,\bm\phi)\) and fixed parameters, but not on old velocity. Since \(\gamma\neq0\), \cref{eq:inverse-velocity} determines velocity; the drift then determines position. Direct substitution verifies both compositions.
\end{proof}

\begin{figure}[H]
\centering
\begin{tikzpicture}[
  box/.style={rounded corners=1.5pt,draw=fiink,fill=fipale,
              minimum width=25mm,minimum height=9mm,align=center},
  known/.style={rounded corners=1.5pt,draw=fiblue,fill=white,
                minimum width=27mm,minimum height=9mm,align=center},
  arr/.style={-{Latex[length=2mm]},line width=0.8pt,color=fiblue},
  node distance=12mm and 16mm
]
\node[box] (rn) {\(\bm r^+\)};
\node[box,right=of rn] (pn) {\(\bm\phi^+\)};
\node[box,right=of pn] (vn) {\(\bm v^+\)};
\node[known,below=of pn] (po) {\(\bm\phi=G_{\bm r^+}^{-1}(\bm\phi^+)\)};
\node[known,below=of vn] (vo) {\(\bm v=\gamma^{-1}[\bm v^+-\dt M^{-1}\bm F]\)};
\node[known,below=of rn] (ro) {\(\bm r=\bm r^+-\bm v\dt\)};
\draw[arr] (rn) -- (po);
\draw[arr] (pn) -- (po);
\draw[arr] (po) -- (vo);
\coordinate (voRoute) at ([xshift=4mm]vo.east);
\coordinate (rnRoute) at ([yshift=4mm]rn.north);
\draw[arr] (rn.north) -- (rnRoute) -- (voRoute |- rnRoute) -- (voRoute) -- (vo.east);
\draw[arr] (vn) -- (vo);
\coordinate (voBottomRoute) at ([yshift=-4mm]vo.south);
\coordinate (roBottomRoute) at ([yshift=-4mm]ro.south);
\draw[arr] (vo.south) -- (voBottomRoute) -- (roBottomRoute) -- (ro.south);
\draw[arr] (rn) -- (ro);
\end{tikzpicture}
\caption{The inverse is triangular: phase branch first, then force and velocity, then position.}
\label{fig:inverse-dag}
\end{figure}

The mechanical sector is therefore globally algebraic for every velocity-independent force when \(\gamma\neq0\). All branch structure enters through the phase map.

\subsection{Full determinant factorization}

\begin{theorem}[Full Jacobian]
\label{thm:determinant}
At every smooth point of the wall-free map with static \(\bm\theta\),
\begin{equation}
\boxed{
\det D\Psi
=
\gamma^{2N}
\det D_{\bm\phi}G_{\bm r^+}(\bm\phi).
}
\label{eq:determinant}
\end{equation}
The determinant is independent of positional and phase derivatives of the force.
\end{theorem}

\begin{proof}
First change coordinates from \((\bm r,\bm v,\bm\phi)\) to
\((\bm r^+=\bm r+\dt\bm v,\bm v,\bm\phi)\); this block-triangular map has determinant one. In these coordinates, the remaining Jacobian is
\[
\begin{pmatrix}
I & 0 & 0\\
\ast & \gamma I & \dt M^{-1}D_\phi\bm F\\
\ast & 0 & D_\phi G_{\bm r^+}
\end{pmatrix}.
\]
Multiplying the diagonal block determinants yields \cref{eq:determinant}.
\end{proof}

\begin{corollary}[Local singularity channels]
The smooth full map is locally invertible exactly where
\[
\gamma\neq0
\qquad\text{and}\qquad
\det D_\phi G_{\bm r^+}\neq0.
\]
Force curvature cannot create a local singularity in this update ordering.
\end{corollary}

\begin{figure}[H]
\centering
\begin{tikzpicture}
\begin{axis}[
  width=0.68\linewidth,
  height=0.45\linewidth,
  xlabel={factorized prediction
    \(\gamma^{2N}\det D_{\phi}G_{\bm r^+}\)},
  ylabel={finite-difference \(\det D\Psi\)},
  xmin=0.18,xmax=1.0,
  ymin=0.18,ymax=1.0,
  axis lines=left,
  grid=both,
  major grid style={draw=firule},
  minor grid style={draw=firule!45},
  tick label style={font=\footnotesize},
  label style={font=\small},
  legend style={draw=none,fill=none,font=\footnotesize,at={(0.03,0.97)},anchor=north west}
]
\addplot[domain=0.18:1.0,figray,dashed,line width=0.8pt] {x};
\addlegendentry{identity line}
\addplot[
  only marks,
  mark=*,
  mark size=2.1pt,
  fiblue
] table[
  x=det_factorized,
  y=det_full,
  col sep=comma
] {evidence/determinant_factorization.csv};
\addlegendentry{15 live-map probes}
\end{axis}
\end{tikzpicture}
\caption{Numerical audit of \cref{eq:determinant}. Maximum relative discrepancy is \(2.1\times10^{-8}\).}
\label{fig:determinant}
\figsource{\sourcecode{evidence/thesis\_verification.json}, \sourcecode{determinant\_factorization}.}
\end{figure}

\subsection{Phase uniqueness}

The phase Jacobian is
\begin{equation}
\begin{aligned}
(D_\phi G)_{ii}
&=
1-\beta\dt
\sum_{j\neq i}
\frac{\cos(\phi_j-\phi_i)}
     {\widetilde d_{ij}^{+}},\\
(D_\phi G)_{ij}
&=
\beta\dt
\frac{\cos(\phi_j-\phi_i)}
     {\widetilde d_{ij}^{+}},
\qquad i\neq j,
\end{aligned}
\label{eq:phase-jacobian}
\end{equation}
where
\(\widetilde d_{ij}^{+}=\norm{\bm r_i^+-\bm r_j^+}+\varepsilon\).

\begin{theorem}[Configuration-dependent contraction certificate]
\label{thm:phase-certificate}
Define
\begin{equation}
q(\bm r^+)
:=
2\abs{\beta}\dt
\max_i\sum_{j\neq i}
\frac{1}{\widetilde d_{ij}^{+}}.
\label{eq:q}
\end{equation}
If \(q(\bm r^+)<1\), then the inverse phase iteration is a contraction in the \(\ell_\infty\) norm on a consistent torus lift. The phase preimage is unique in that lift, fixed-point inversion converges geometrically, and \(G_{\bm r^+}\) has lower bi-Lipschitz bound \(1-q\).
\end{theorem}

The certificate is configuration-dependent. It supplies an operational gate for every trajectory used in the round-trip study.

\subsection{Global branch geometry}

For two tokens at fixed softened separation \(D\), let
\[
\delta=\phi_1-\phi_2,
\qquad
\Delta\omega=\omega_1-\omega_2,
\qquad
a=\frac{2\beta\dt}{D}.
\]
The exact relative-phase map is
\begin{equation}
\boxed{
\delta^+
=
\delta+\Delta\omega\,\dt-a\sin\delta
\pmod{2\pi}.
}
\label{eq:two-clock}
\end{equation}
For equal frequencies, its derivative is \(1-a\cos\delta\). The circle map is a diffeomorphism for \(\abs a<1\), a singular monotone homeomorphism at \(\abs a=1\), and noninjective for \(\abs a>1\).

\begin{figure}[H]
\centering
\begin{tikzpicture}
\begin{axis}[
  width=0.72\linewidth,
  height=0.43\linewidth,
  xlabel={input phase difference \(\delta\)},
  ylabel={lifted output \(f_a(\delta)=\delta-a\sin\delta\)},
  xmin=-3.2,xmax=3.2,
  ymin=-4.2,ymax=4.2,
  grid=major,
  major grid style={draw=firule},
  tick label style={font=\footnotesize},
  label style={font=\small},
  legend style={draw=none,fill=none,font=\footnotesize,at={(0.03,0.97)},anchor=north west}
]
\addplot[fiblue,line width=1.2pt,domain=-3.1416:3.1416,samples=180]
  {x-0.6*sin(deg(x))};
\addlegendentry{\(a=0.6\): one branch}
\addplot[fiorange,line width=1.2pt,domain=-3.1416:3.1416,samples=180]
  {x-1.6*sin(deg(x))};
\addlegendentry{\(a=1.6\): folded branches}
\addplot[figray,dashed,domain=-3.1416:3.1416] {x};
\end{axis}
\end{tikzpicture}
\caption{The same phase law contains an injective regime and a folded regime. Distinction provenance locates the transition between them.}
\label{fig:phase-fold}
\end{figure}

The live eight-token system reaches the folded regime at default parameters. At seed 22, step 9, one target has \(q=1.338984\) and three refined phase roots. Combining each root with the algebraic mechanical inverse produces three complete pre-states.

\begin{table}[H]
\centering
\caption{Three complete pre-states for one target. Every phase Jacobian is locally nonsingular; the distinction is global branch structure.}
\label{tab:multiroot}
\small
\begin{tabular}{@{}rrrr@{}}
\toprule
root & \(\det D_\phi G\) &
distance from generating pre-state &
forward residual\\
\midrule
0 & \(-0.2300\) & \(1.09\times10^{-15}\) & \(0\)\\
1 & \( 0.4584\) & \(7.386\) & \(1.50\times10^{-13}\)\\
2 & \( 0.4115\) & \(6.738\) & \(1.33\times10^{-13}\)\\
\bottomrule
\end{tabular}
\figsource{\sourcecode{evidence/thesis\_verification.json}, \sourcecode{phase\_multiroot\_collision}.}
\end{table}

This example demonstrates the value of the framework: a nonzero local Jacobian, a convergent inverse solver, and global branch uniqueness are separate questions.

\subsection{Finite analogue}

On \(\Zmod/M\Zmod\), consider
\begin{equation}
r^+=r+v\pmod M,
\qquad
v^+=g\,v+F(r^+)\pmod M.
\label{eq:finite-map}
\end{equation}

\begin{theorem}[Modular unit criterion]
\label{thm:finite}
\Cref{eq:finite-map} is a permutation of the \(M^2\) states for every force \(F\) if and only if \(\gcd(g,M)=1\).
\end{theorem}

\begin{proof}
If \(g\) is a unit,
\[
v=g^{-1}[v^+-F(r^+)]\pmod M,
\qquad
r=r^+-v\pmod M.
\]
If \(g\) is not a unit, multiplication by \(g\) has a nontrivial kernel, and compensated position--velocity pairs collide.
\end{proof}

Enumeration gives \(63{,}001/63{,}001\) images for \((M,g)=(251,3)\), while \((M,g)=(16,2)\) yields \(128/256\) images. The theorem provides a complete finite model of the distinction between invertible multiplication and state aggregation.

\begin{remark}[Six independent properties]
Bijectivity, local nonsingularity, time-reversal symmetry, volume preservation, energy conservation, and stable numerical inversion are distinct properties. The results above identify exactly which of them follow from each assumption.
\end{remark}
